\abstract{Wildfires are a serious problem in Kazakhstan. Every year, satellites detect thousands of new vegetation fires, and in June 2023 a fire in the Abay Region killed 14 forestry workers. A daily map that shows where new fires are likely could help emergency services prepare. This study aimed to build and test a machine-learning model that predicts, one day ahead, the chance of a new fire in every 5\,km square of Kazakhstan and explains what drives this chance. We combined satellite fire detections for 2015--2024 with weather, vegetation, terrain, land-cover, and population data. Four machine-learning models were trained on 2015--2021, calibrated on 2022, and tested on 2023--2024. They were compared with two simple references: past fire frequency by season and a weather-only model. The best model, XGBoost, ranked the squares about 104 times better than a random ranking. The 5\% of the area with the highest daily scores contained 55.6\% of all new fires, compared with 14.3\% for the seasonal reference. Terrain, land cover, population, and season were the most useful factors, while week-old weather added little. The model worked well in later years but much worse in regions it had not seen, and its probabilities were 1.2--1.8 times too high. Therefore, the model is most useful as a daily ranking tool, not as an exact probability. It can show emergency services where to look first, and it shows that fire models should be tested at the real fire rate.}

\keywords{Wildfire prediction, Kazakhstan, Machine learning, XGBoost, Satellite remote sensing, Probability calibration, SHAP}
